\subsection{A Non-Linear System} 
The graphs of linear systems are \makeblank. This gives a major an advantage in solving these systems: two lines in a plane intersect \makeblank[1in] most of the time, or \makeblank[1.5in] if they're parallel. What about nonlinear equations?
\begin{Example} Consider the system:
\begin{syspic}[0]{4}
\draw (-.5\textwidth,0) node{$
\begin{systemL}
x^2+y^2=5&\\
y=x^2-3&
\end{systemL}$};
\draw[graphend] (0,0) circle ({sqrt(5)});
\draw[graph,domain={-sqrt(7)}:{sqrt(7)},samples=100] plot (\x,\x*\x-3);
\end{syspic}
\begin{itemize}
\item The graph of $(I)$ is a \makeblank[1in] and the graph of $(II)$ is a \makeblank[1in]
\item The curves intersect at \makeblanksmall points!
\item Each point is a solution of the system:
\end{itemize}
\begin{lpic}{}{3}
\end{lpic}
\end{Example}
\begin{Example} Consider the system:
\begin{syspic}[0]{4}
\draw (-.5\textwidth,0) node{$
\begin{systemL}
(x-1)^2+(y-1)^2=2&\\
x^2+y^2=9&
\end{systemL}$};
\draw[graphend] (1,1) circle ({sqrt(2)});
\draw[graphend] (0,0) circle (3);
\end{syspic}
\begin{itemize}
\item The graphs of these equations are both \makeblank[1in].
\item These curves \makeblank.
\item This system has \makeblank.
\end{itemize}
\end{Example}
Systems of nonlinear systems might have:
\begin{itemize}
\item No solution because the curves go around each other
\item Several solutions because the curves cross each other more than once
\item All the complications involved in solving nonlinear equations
\end{itemize}